\appendix
\section{\texorpdfstring{Worked Example: $2$-dimensional system $\otimes$ $2$-dimensional environment}{Worked Example: 2D system x 2D environment}}
\label{app:worked-example}

This appendix gives a fully explicit instance of an unfaithful cut and shows how the associated
$\Sigma_2$ object is constructed.

\subsection{Setup}

Let $\mathcal{H}$ be the set of admissible histories of a joint system $SE$ with Hilbert space
$\mathcal{H}_S \otimes \mathcal{H}_E$, where $\dim \mathcal{H}_S=\dim \mathcal{H}_E=2$.
Let $\Pi$ be the reduction to the system density matrix at a chosen reduction event:
\[
  \Pi(h) = \rho_S(t_0) = \mathrm{Tr}_E\big[\rho_{SE}(t_0)\big].
\]
Let the intervention class $\mathcal{I}$ include free evolution for a fixed duration $\Delta t$,
implemented by a joint unitary $U_{SE}$:
\[
  I_{\Delta t}: \rho_{SE}(t_0) \mapsto \rho_{SE}(t_0+\Delta t) = U_{SE}\rho_{SE}(t_0)U_{SE}^\dagger,
\quad
  \rho_S(t_0+\Delta t)=\mathrm{Tr}_E[\rho_{SE}(t_0+\Delta t)].
\]
Define the intervention relation $h \sim_{I_{\Delta t}} h'$ to mean that $h'$ is a possible
post-intervention continuation of $h$ under $I_{\Delta t}$.

\subsection{Two distinct histories with identical reduced state}

Fix an orthonormal basis $\{|0\rangle_S,|1\rangle_S\}$ for $S$ and two distinct orthonormal
bases $\{|0\rangle_E,|1\rangle_E\}$ and $\{|0'\rangle_E,|1'\rangle_E\}$ for $E$.
Define two pure joint states at $t_0$:
\begin{align}
  |\psi_1\rangle_{SE} &= \sqrt{p}\,|0\rangle_S|0\rangle_E + \sqrt{1-p}\,|1\rangle_S|1\rangle_E,\\
  |\psi_2\rangle_{SE} &= \sqrt{p}\,|0\rangle_S|0'\rangle_E + \sqrt{1-p}\,|1\rangle_S|1'\rangle_E,
\end{align}
with $p\in(0,1)$ and with $\{|0\rangle_E,|1\rangle_E\}$ not equal to $\{|0'\rangle_E,|1'\rangle_E\}$.

Both reduce to the same system state:
\[
  \rho_S(t_0)
  = \mathrm{Tr}_E\big[|\psi_1\rangle\langle\psi_1|\big]
  = \mathrm{Tr}_E\big[|\psi_2\rangle\langle\psi_2|\big]
  = p\,|0\rangle\langle 0| + (1-p)\,|1\rangle\langle 1|.
\]
Thus, there exist distinct histories $h_1,h_2 \in \mathcal{H}$ with $\Pi(h_1)=\Pi(h_2)$.

\subsection{Explicit correlation-sensitive dynamics}

Choose an entangling unitary, e.g.\ controlled-NOT with $S$ as control and $E$ as target, in the
$\{|0\rangle_E,|1\rangle_E\}$ basis:
\[
  U_{SE} = |0\rangle\langle 0|_S \otimes \mathbb{I}_E + |1\rangle\langle 1|_S \otimes X_E,
\]
where $X_E|0\rangle_E=|1\rangle_E$ and $X_E|1\rangle_E=|0\rangle_E$.

Under this unitary, $|\psi_1\rangle_{SE}$ evolves to
\[
  |\psi_1'\rangle_{SE} = \sqrt{p}\,|0\rangle_S|0\rangle_E + \sqrt{1-p}\,|1\rangle_S X_E|1\rangle_E
  = \sqrt{p}\,|0\rangle_S|0\rangle_E + \sqrt{1-p}\,|1\rangle_S|0\rangle_E,
\]
which factorizes as $(\sqrt{p}\,|0\rangle_S+\sqrt{1-p}\,|1\rangle_S)\otimes|0\rangle_E$ and hence
yields a pure reduced state
\[
  \rho_S^{(1)}(t_0+\Delta t) = |\phi\rangle\langle\phi|,\quad |\phi\rangle=\sqrt{p}\,|0\rangle+\sqrt{1-p}\,|1\rangle.
\]

By contrast, if $|\psi_2\rangle_{SE}$ is defined using a different environment basis, then the
same $U_{SE}$ (defined relative to the unprimed basis) acts differently on the environment
components of $|\psi_2\rangle_{SE}$. For generic choices of $\{|0'\rangle,|1'\rangle\}$, the
resulting reduced state $\rho_S^{(2)}(t_0+\Delta t)$ differs from $|\phi\rangle\langle\phi|$.
Therefore,
\[
  \Pi(h_1)=\Pi(h_2)\quad\text{but}\quad
  \Pi(h_1') \neq \Pi(h_2')
\]
for $h_1 \sim_{I_{\Delta t}} h_1'$ and $h_2 \sim_{I_{\Delta t}} h_2'$.

This is a predictive failure: an agent who knows only $\rho_S(t_0)$ cannot predict the outcome
of applying the same free-evolution intervention $I_{\Delta t}$, because the result depends on
which fiber element (which correlation structure) is realized.

\subsection{Constructing the \texorpdfstring{$\Sigma_2$}{Sigma-2} object}

Let $r=\rho_S(t_0)$. The fiber is $F(r)=\Pi^{-1}(r)$, containing at least $h_1,h_2$.
For an intervention $I_{\Delta t}$, define the continuation set
\[
  A_{I_{\Delta t}}(r)
  = \Big\{ \Pi(h') \,\Big|\, \exists h\in \Pi^{-1}(r)\ \text{s.t.}\ h \sim_{I_{\Delta t}} h' \Big\}.
\]
In this example, $A_{I_{\Delta t}}(r)$ contains at least two distinct elements
$\rho_S^{(1)}(t_0+\Delta t)$ and $\rho_S^{(2)}(t_0+\Delta t)$, reflecting the set-valued nature of
faithful prediction under the unfaithful cut.

The associated $\Sigma_2$ description is
\[
  \Sigma_2(r) = \big(r,\ F(r),\ \{A_I(r)\}_{I\in\mathcal{I}}\big),
\]
which restores counterfactual semantics: given $I\in\mathcal{I}$, the correct prediction is not
a single output but the admissible outcome set determined by the fiber and the intervention
relation.